%% La date de Pâques la plus tôt est le 25 mars
%% La date de Pentecôte la plus tard est le 13 juin
%% => sanctoral du 25 mars (transféré) au 12 juin, moins les fêtes non transférées qui tombent au pire dans l'octave de Pâques (27 28 mars)

\festum{0325}{}{25 mars}
\saintindex{}
\saintindex{}
\oratio{Deus, qui de beátæ Maríæ Vírginis útero Verbum tuum, Angelo nuntiánte, carnem suscípere voluísti: præsta supplícibus tuis; ut, qui vere eam Genetrícem Dei crédimus, ejus apud te intercessiónibus adjuvémur. \pereumdemla}{Ô Dieu, qui avez voulu que votre Verbe prît un corps humain à la parole de l’Ange dans le sein de la bienheureuse Vierge Marie ; accordez à ceux qui vous en supplient que, nous qui la croyons véritablement Mère de Dieu, nous soyons secourus auprès de vous grâce à son intercession. \pereumdemfr}
\festum{0401}{}{1 avril}
\saintindex{}
\saintindex{}
\oratio{Deus omípotens, cujus sérvitus summa et plena est felícitas: * concéde nobis, intercedénte beáto Hugóne, Abbáte, § in sanctitáte et justítia sine timóre tibi servíre. \perdominumla}{}
\festum{0402}{S. Francisci de Paula Confessoris}{2 avril}
\saintindex{}
\saintindex{}
\oratio{Deus, humílium celsitúdo, qui beátum Francíscum Confessórem Sanctórum tuórum glória sublimásti: tríbue, quǽsumus; ut, ejus méritis et imitatióne, promíssa humílibus prǽmia felíciter consequámur. \perdominumla.}{Ô Dieu, grandeur des humbles, qui avez élevé le bienheureux François, Confesseur, à la gloire de vos Saints, accordez, nous vous en prions, que, par ses mérites et par l’imitation de ses vertus, nous ayons le bonheur d’obtenir les récompenses promises à ceux qui sont humbles. \perdominumfr.}
\festum{0404}{S. Isidori Episcopi Confessoris et Ecclesiæ Doctoris}{4 avril}
\saintindex{Isidórum}
\saintindex{Isidóre}
\oratio{}{Ô Dieu qui avez fait à votre peuple la grâce d’avoir le bienheureux Isidore, pour ministre du salut éternel, faites, nous vous en prions, que nous méritions d’avoir pour intercesseur dans les cieux celui qui nous a donné sur terre la doctrine de vie. \perdominumfr}
\festum{0405}{S. Vincentii Ferrerii Confessoris}{5 avril}
\saintindex{}
\saintindex{}
\oratio{Deus, qui Ecclésiam tuam beáti Vincéntii Confessóris tui méritis et prædicatióne illustráre dignátus es: concéde nobis fámulis tuis; ut et ipsíus instruámur exémplis et ab ómnibus ejus patrocínio liberémur advérsis. \perdominumla}{Dieu, qui avez daigné illustrer Votre Église par Les mérites et la prédication du bienheureux Vincent, Votre confesseur, accordez à nous, Vos familiers, que nous soyons instruits par ses exemples et libérés de toutes adversités par son patronage. \perdominumfr}
\festum{0411}{S. Leonis I Papæ Confessoris et Ecclesiæ Doctoris}{11 avril}
\saintindex{Leónem}
\saintindex{Léon}
\oratio{@Commune/C4-1}{Nous vous supplions, Seigneur, d’exaucer les prières que nous vous adressons en la solennité du bienheureux Léon, votre Confesseur et Pontife, et de nous accorder, grâce aux mérites et à l’intercession de celui qui vous a si dignement servi, le pardon de tous nos péchés. \perdominumfr}
\festum{0413}{S. Hermenegildi Martyris}{13 avril}
\saintindex{}
\saintindex{}
\oratio{Deus, qui beátum Hermenegíldum Mártyrem tuum cælésti regno terrénum postpónere docuísti: da quǽsumus nobis; ejus exémplo cadúca despícere, atque ætérna sectári. \perdominumla}{Ô Dieu, qui avez appris au Bienheureux Herménégilde à faire passer le royaume de ce monde après celui des cieux, accordez-nous, s’il Vous plaît, de mépriser, à son exemple, les biens périssables et de poursuivre les éternels. \perdominumfr}
\festum{0414}{S. Justini Martyris}{14 avril}
\saintindex{}
\saintindex{}
\oratio{Deus, qui per stultítiam crucis eminéntem Jesu Christi sciéntiam beátum Justínum Mártyrem mirabíliter docuísti: ejus nobis intercessióne concéde; ut, errórum circumventióne depúlsa, fídei firmitátem consequámur. \pereumdemla}{Ô Dieu qui, par la folie de la croix, avez merveilleusement enseigné au bienheureux Martyr Justin la suprême science de Jésus-Christ, accordez-nous par Son intercession que, repoussant tout piège d’erreur, nous obtenions la fermeté de la foi. \pereumdemfr}
\festum{0414}{Ss. Tiburtii, Valeriani, et Maximi Martyrum}{14 avril}
\saintindex{}
\saintindex{}
\oratio{Præsta, quǽsumus, omnípotens Deus: ut, qui sanctórum Mártyrum tuórum Tiburtii, Valeriáni et Maximi solemnia colimus; eórum étiam virtútes imitemur. \perdominumla}{Faites, nous Vous en supplions, Seigneur, que nous qui célébrons la fête de Vos saints Martyrs Tiburce, Valérien et Maxime, nous imitions aussi leurs vertus. \perdominumfr}
\festum{0417}{S. Aniceti Papæ et Martyris}{17 avril}
\saintindex{Anicéti}
\saintindex{Anicet}
\festum{0421}{S. Anselmi Episcopi Confessoris et Ecclesiæ Doctoris}{21 avril}
\saintindex{Ansélmum}
\saintindex{Anselme}
\festum{0422}{SS. Soteris et Caji Summorum Pontificum et Martyrum}{22 avril}
\saintindex{Sotéris et Caji}
\saintindex{Soter et Caius}
\festum{0423}{S. Georgii Martyris}{23 avril}
\saintindex{}
\saintindex{}
\oratio{Deus, qui nos beáti Geórgii Mártyris tui méritis et intercessióne lætíficas: concéde propítius; ut, qui tua per eum benefícia póscimus, dono tuæ grátiæ consequámur. \perdominumla}{Ô Dieu qui nous réjouissez par les mérites et l’intercession du Bienheureux Georges, Votre Martyr, accordez-nous, dans Votre bonté, que, Vous demandant par lui Vos bienfaits, nous les obtenions par le don de Votre grâce. \perdominumfr}
\festum{0424}{S. Fidelis de Sigmaringa Martyris}{24 avril}
\saintindex{}
\saintindex{}
\oratio{Deus, qui beátum Fidélem, seráphico spíritus ardóre succénsum, in veræ fídei propagatióne martýrii palma et gloriósis miráculis decoráre dignátus es: ejus, quǽsumus, méritis et intercessióne, ita nos per grátiam tuam in fide et caritáte confírma; ut in servítio tuo fidéles usque ad mortem inveníri mereámur. \perdominumla}{Ô Dieu qui, après avoir embrasé le Bienheureux Fidèle de l’ardeur de l’esprit séraphique, avez daigné l’illustrer par la palme du martyre et par de glorieux miracles, nous Vous demandons, par ses mérites et son intercession, de si bien nous confirmer, par Votre grâce, dans la foi et la charité, que nous méritions d’être trouvés fidèles dans Votre service jusqu’à la mort. \perdominumfr}
\festum{0425}{S. Marci Evangelistæ}{25 avril}
\saintindex{}
\saintindex{}
\oratio{Deus, qui beátum Marcum Evangelístam tuum evangélicæ prædicatiónis grátia sublimásti: tríbue, quǽsumus; ejus nos semper et eruditióne profícere et oratióne deféndi. \perdominumla}{Ô Dieu, qui avez élevé si haut le bienheureux Marc, Votre Evangéliste, par la grâce de la prédication évangélique, accordez-nous, s’il Vous plaît, de profiter toujours de sa doctrine et d’être toujours défendus par sa prière. \perdominumfr}
\festum{0426}{SS. Cleti et Marcellini Summorum Pontificum et Martyrum}{26 avril}
\saintindex{Cleti et Marcellíni}
\saintindex{Clet et Marcellin}
\festum{0427}{S. Petri Canisii Confessoris et Ecclesiæ Doctoris}{27 avril}
\saintindex{Petrum}
\saintindex{Pierre}
\oratio{Deus, qui ad tuéndam cathólicam fidem beátum Petrum Confessórem tuum virtúte et doctrína roborásti: concéde propítius; ut ejus exémplis et mónitis errántes ad salútem resipíscant, et fidéles in veritátis confessióne persevérent. \perdominumla}{Ô Dieu, qui pour la protection de la foi catholique, avez affermi dans la vertu et la doctrine le Bienheureux Pierre, Votre Confesseur, accordez, dans Votre clémence, qu’instruits par ses exemples et ses enseignements, les égarés reviennent à la voie du salut et que les fidèles persévèrent dans la confession de la vérité. \perdominumfr}
\festum{0428}{S. Pauli a Cruce Confessoris}{28 avril}
\saintindex{}
\saintindex{}
\oratio{Dómine Jesu Christe, qui, ad mystérium crucis prædicándum, sanctum Paulum singulári caritáte donásti, et per eum novam in Ecclésia famíliam floréscere voluísti: ipsíus nobis intercessióne concéde; ut, passiónem tuam júgiter recoléntes in terris, ejúsdem fructum cónsequi mereámur in cælis:
$Qui vivis}{Seigneur Jésus-Christ qui, pour la prédication du mystère de la Croix, avez gratifié saint Paul d’une charité sans pareille et qui avez voulu, par lui, faire fleurir une nouvelle famille dans l’Eglise, accordez-nous, par son intercession, qu’en nous remémorant continuellement Votre passion sur Terre, nous méritions d’en recueillir les fruits au ciel,
$Qui vivis}
\festum{0428}{S. Peter Chanel, priest and martyr}{28 avril}
\saintindex{Petri}
\saintindex{Pierre}
\festum{0428}{S. Vitalis Martyris}{28 avril}
\saintindex{Vitális}
\saintindex{Vital}
\festum{0429}{S. Petri Martyris}{29 avril}
\saintindex{}
\saintindex{}
\oratio{Præsta, quǽsumus, omnípotens Deus: ut beáti Petri Mártyris tui fidem cóngrua devotióne sectémur; qui, pro ejúsdem fídei dilatatióne, martýrii palmam méruit obtinére. \perdominumla}{Faites, nous Vous le demandons, ô Dieu tout-puissant, que nous imitions avec le zèle qui convient, la foi du bienheureux Pierre, Votre Martyr qui, pour l’expansion de cette même foi, a mérité d’obtenir la palme du martyre. \perdominumfr}
\festum{0430}{S. Catharinæ Senensis Virginis}{30 avril}
\saintindex{}
\saintindex{}
\oratio{Da, quǽsumus, omnípotens Deus: ut, qui beátæ Catharínæ Vírginis tuæ natalítia cólimus; et ánnua solemnitáte lætémur, et tantæ virtútis proficiámus exémplo. \perdominumla}{Accordez-nous, s’il Vous plaît, Dieu tout-puissant, qu’honorant la naissance au ciel de Votre bienheureuse Vierge Catherine, nous nous réjouissions de cette solennité annuelle et progressions par l’exemple d’une si grande vertu. \perdominumfr}
\festum{0501}{Ss. Philippi et Jacobi Apostolorum}{1 mai}
\saintindex{}
\saintindex{}
\oratio{Deus, qui nos ánnua Apostolórum tuórum Philíppi et Jacóbi solemnitáte lætíficas: præsta, quǽsumus; ut, quorum gaudémus méritis, instruámur exémplis. \perdominumla}{Dieu qui nous réjouissez en la solennité annuelle de Vos Apôtres Philippe et Jacques ; faites, nous Vous en prions, que nous soyons instruits aux exemples de ceux dont les mérites nous remplissent d’allégresse. \perdominumfr}
\festum{0501}{S. Joseph Sponsi B.M.V. Confessoris}{1 mai}
\saintindex{}
\saintindex{}
\oratio{Deus, qui ineffábili providéntia beátum Joseph sanctíssimæ Genitrícis tuæ sponsum eligere dignátus es: præsta, quǽsumus; ut, quem protectórem venerámur in terris, intercessórem habere mereámur in cælis:
$Qui vivis}{Dieu créateur de toutes choses, qui avez imposé au genre humain la loi du travail : accordez-nous favorablement ; grâce à l’exemple et au patronage de saint Joseph, d’accomplir parfaitement le travail que vous nous fixez, et d’obtenir les récompenses que vous nous promettez.
$Qui vivis}
\festum{0501}{S. Joseph Opificis}{1 mai}
\saintindex{}
\saintindex{}
\oratio{Rerum cónditor Deus, qui legem labóris humáno géneri statuísti: concéde propítius; ut, sancti Joseph exémplo et patrocínio, ópera perficiámus quæ prǽcipis, et prǽmia consequámur quæ promíttis. \perdominumla}{Créateur du monde, ô Dieu, qui avez imposé la loi du travail au genre humain, accordez-nous avec bonté que, à l’exemple et sous le patronage de saint Joseph, nous accomplissions les travaux que vous ordonnez et nous obtenions les récompenses que vous promettez. \perdominumfr}
\festum{0502}{S. Athanasii Episcopi Confessoris et Ecclesiæ Doctoris}{2 mai}
\saintindex{Athanásii}
\saintindex{Athanase}
\oratio{@Commune/C4-1}{}
\festum{0503}{}{3 mai}
\saintindex{}
\saintindex{}
\oratio{Deus, qui in præclára salutíferæ Crucis Inventióne, passiónis tuæ mirácula suscitásti: concéde; ut vitális ligni prétio, ætérnæ vitæ suffrágia consequámur:
$Qui vivis}{Dieu, qui lors de la glorieuse Découverte de la Croix, instrument de notre salut, avez renouvelé les miracles de votre passion : faites qu’au prix de cet arbre de vie, nous méritions d’obtenir la vie éternelle.
$Qui vivis}
\festum{0503}{Ss. Alexandri I Papæ, Eventii et Theoduli Mm. ac Juvenalis Ep. et Conf.}{3 mai}
\saintindex{}
\saintindex{}
\oratio{Præsta, quǽsumus, omnípotens Deus: ut, qui sanctórum tuórum Alexándri, Evéntii, Theodúli atque Juvenális natalítia cólimus; a cunctis malis imminéntibus, eórum intercessiónibus, liberémur.  \perdominumla}{Accordez-nous, s’il vous plaît, ô Dieu tout-puissant, que, célébrant la naissance au ciel de vos saints Alexandre, Eventius, Théodule et Juvénal, nous soyons délivrés, grâce à leur intercession, de tous les maux qui nous menacent. \perdominumfr}
\festum{0503}{}{3 mai}
\saintindex{}
\saintindex{}
\festum{0504}{S. Monicæ Viduæ}{4 mai}
\saintindex{Mónicæ}
\saintindex{Monique}
\oratio{Deus, mæréntium consolátor et in te sperántium salus, qui beátæ Mónicæ pias lácrimas in conversióne fílii sui Augustíni misericórditer suscepísti: da nobis utriúsque intervéntu; peccáta nostra deploráre, et grátiæ tuæ indulgéntiam inveníre. \perdominumla}{Ô Dieu, consolateur des affligés et salut de ceux qui espèrent en Vous, Vous avez miséricordieusement exaucé les pieuses larmes versées par la bienheureuse Monique pour la conversion de son fils Augustin ; donnez-nous, par l'intercession de tous deux, de pleurer nos péchés et de trouver l'indulgence de Votre grâce. \perdominumfr}
\festum{0505}{S. Pii V Papæ et Confessoris}{5 mai}
\saintindex{Pium}
\saintindex{Pie}
\oratio{Deus, qui, ad conteréndos Ecclésiæ tuæ hostes et ad divínum cultum reparándum, beátum Pium Pontíficem Máximum elígere dignátus es: fac nos ipsíus deféndi præsídiis et ita tuis inhærére obséquiis; ut, ómnium hóstium superátis insídiis, perpétua pace lætémur. \perdominumla}{Ô Dieu qui, pour écraser les ennemis de Votre Eglise et restaurer le culte divin, avez daigné choisir le bienheureux Pie pour Souverain Pontife, faites que, défendus par son secours pour demeurer très attachés à Votre service, nous triomphions des pièges de tous nos ennemis et jouissions de l'éternelle paix. \perdominumfr}
\festum{0506}{S. Joannis Apostoli ante Portam Latinam}{6 mai}
\saintindex{}
\saintindex{}
\oratio{Deus, qui conspicis, quia nos undique mala nostra perturbant: præsta, quǽsumus; ut beáti Joánnis Apóstoli tui et Evangelistæ intercéssio gloriósa nos prótegat. \perdominumla}{Ô Dieu qui voyez comment nous sommes troublés par les maux qui nous arrivent de partout, faites, nous vous le demandons, que nous soyons protégés par la glorieuse intercession du bienheureux Jean, votre Apôtre et Évangéliste.  \perdominumfr}
\festum{0507}{S. Stanislai Episcopi et Martyris}{7 mai}
\saintindex{Stanisláus}
\saintindex{Stanislas}
\oratio{Deus, pro cujus honóre gloriósus Póntifex Stanisláus gládiis impiórum occúbuit: præsta, quǽsumus; ut omnes, qui ejus implórant auxílium, petitiónis suæ salutárem consequántur efféctum. \perdominumla}{Ô Dieu, c'est pour Votre honneur que le glorieux Pontife Stanislas a succombé sous le glaive des impies ; accordez-nous, s'il Vous plaît, que tous ceux qui implorent son aide, obtiennent pour leur salut l'effet de leur demande. \perdominumfr}
\festum{0508}{In Apparitione S. Michaëlis Archangeli}{8 mai}
\saintindex{}
\saintindex{}
\oratio{Deus, qui, miro órdine, Angelórum ministéria hominúmque dispénsas: concéde propítius; ut, a quibus tibi ministrántibus in cælo semper assístitur, ab his in terra vita nostra muniátur.  \perdominumla}{Ô Dieu, qui distribuez selon un ordre admirable les ministères des Anges et des hommes, accordez-nous miséricordieusement que ceux qui, dans le ciel, Vous entourent d'un continuel service, soient la protection de notre vie sur terre. \perdominumfr}
\festum{0509}{S. Gregorii Nazianzeni Episcopi Confessoris et Ecclesiæ Doctoris}{9 mai}
\saintindex{Gregórium}
\saintindex{Grégoire }
\festum{0509}{}{9 mai}
\saintindex{}
\saintindex{}
\festum{0510}{S. Antonini Episcopi et Confessoris}{10 mai}
\saintindex{}
\saintindex{}
\oratio{Sancti Antoníni, Dómine, Confessóris tui atque Pontíficis méritis adjuvémur: ut, sicut te in illo mirábilem prædicámus, ita in nos misericórdem fuísse gloriémur. \perdominumla}{Que les mérites de saint Antonin, votre Confesseur et Pontife, nous soient en aide, ô Seigneur ; et comme nous vous proclamons admirable dans votre serviteur, faites que nous puissions aussi nous glorifier de votre miséricorde à notre égard. \perdominumfr}
\festum{0510}{Ss. Gordiani et Epimachi Martyrum}{10 mai}
\saintindex{}
\saintindex{}
\oratio{Da, quǽsumus, omnípotens Deus: ut, qui beatórum Mártyrum tuórum Gordiáni et Epímachi solémnia cólimus, eórum apud te intercessiónibus adjuvémur. \perdominumla}{Faites, s’il vous plaît, Dieu tout-puissant, qu’en célébrant la fête de vos bienheureux martyrs Gordien et Épimaque, nous soyons secourus par leur intercession auprès de vous.  \perdominumfr}
\festum{0511}{}{11 mai}
\saintindex{}
\saintindex{}
\festum{0512}{Ss. Nerei, Achillei et Domitillæ Virg. atque Pancratii Martyrum}{12 mai}
\saintindex{}
\saintindex{}
\oratio{Semper nos, Dómine, Mártyrum tuórum Nérei, Achíllei, Domitíllæ atque Pancrátii fóveat, quǽsumus, beáta solémnitas: et tuo dignos reddat obséquio. \perdominumla}{Que toujours, nous Vous en supplions, Seigneur, l’heureuse fête de Vos Martyrs Nérée, Achillée, Domitille et Pancrace nous soutienne et nous rende dignes de Votre service. \perdominumfr}
\festum{0513}{S. Roberti Bellarmino Episcopi Confessoris et Ecclesiæ Doctoris}{13 mai}
\saintindex{Robérte}
\saintindex{Robert}
\oratio{Deus, qui ad errórum insídias repelléndas et apostólicæ Sedis jura propugnánda, beátum Robértum Pontíficem tuum atque Doctórem mira eruditióne et virtúte decorásti: ejus méritis et intercessióne concéde; ut nos in veritátis amóre crescámus et errántium corda ad Ecclésiæ tuæ rédeant unitátem. \perdominumla}{Ô Dieu, pour repousser les pièges de l’erreur et défendre les droits du Siège apostolique, Vous avez doté le bienheureux Robert, Votre Pontife et Docteur, d’une science et d’une force admirable ; faites, par ses mérites et son intercession, que nous grandissions dans l’amour de la vérité et que les cœurs des égarés reviennent à l’unité de Votre Église. \perdominumfr}
\festum{0513}{}{13 mai}
\saintindex{}
\saintindex{}
\festum{0514}{S. Bonifatii Martyris}{14 mai}
\saintindex{}
\saintindex{}
\oratio{Da, quǽsumus, omnípotens Deus: ut qui beáti Bonifátii Mártyris tui solémnia cólimus, ejus apud te intercessiónibus adjuvémur. \perdominumla}{Faites, nous Vous en prions, Dieu tout-puissant, que nous, qui célébrons la fête de Votre bienheureux Martyr Boniface, nous recevions l'assistance de son intercession. \perdominumfr}
\festum{0515}{S. Joannis Baptistæ de la Salle Confessoris}{15 mai}
\saintindex{}
\saintindex{}
\oratio{Deus, qui, ad christiánam páuperum eruditiónem et ad juvéntam in via veritátis firmándam, sanctum Joánnem Baptístam Confessórem excitásti, et novam per eum in Ecclésia famíliam collegísti: concéde propítius; ut ejus intercessióne et exémplo, stúdio glóriæ tuæ in animárum salúte fervéntes, ejus in cælis corónæ partícipes fíeri valeámus. \perdominumla}{Ô Dieu qui, pour l’instruction chrétienne des pauvres et pour la confirmation de la jeunesse dans la voie de la vérité, avez suscité Votre confesseur Jean-Baptiste et avez par lui rassemblé dans l’Église une nouvelle famille, accordez-nous dans Votre bonté, qu’à son exemple et par son intercession, brûlants de zèle pour procurer Votre gloire au moyen du salut des âmes, nous puissions dans les cieux partager sa récompense. \perdominumfr}
\festum{0516}{S. Ubaldi Episcopi et Confessoris}{16 mai}
\saintindex{}
\saintindex{}
\oratio{Auxílium tuum nobis, Dómine, quǽsumus, placátus impénde: et, intercessióne beáti Ubáldi Confessóris tui atque Pontíficis, contra omnes diáboli nequítias déxteram super nos tuæ propitiatiónis exténde. \perdominumla}{Laissez-Vous fléchir, Seigneur : que l’intercession du bienheureux Ubald, Votre Confesseur et Pontife, nous obtienne Votre secours ; étendez sur nous Votre main miséricordieuse pour nous défendre contre toutes les perfidies du démon. \perdominumfr}
\festum{0516}{}{16 mai}
\saintindex{}
\saintindex{}
\oratio{Deus, qui ob invictum beati Joannis sacramentale silentium, nova Ecclesiam tuam martyrii corona decorasti: da nobis, ejus intercessione et exemplo, linguam caute custodire; ac omnia potius mala, quam animæ detrimentum, in hoc sæculo tolerare. \perdominumla.}{Deus, qui ob invictum beati Joannis sacramentale silentium, nova Ecclesiam tuam martyrii corona decorasti: da nobis, ejus intercessione et exemplo, linguam caute custodire; ac omnia potius mala, quam animæ detrimentum, in hoc sæculo tolerare. \perdominumfr.}
\festum{0517}{S. Paschalis Baylon Confessoris}{17 mai}
\saintindex{}
\saintindex{}
\oratio{Deus, qui beátum Paschálem Confessórem tuum mirífica erga Córporis et Sánguinis tui sacra mystéria dilectióne decorásti: concéde propítius; ut, quam ille ex hoc divíno convívio spíritus percépit pinguédinem, eándem et nos percípere mereámur:
$Qui vivis}{Ô Dieu, qui avez orné l’âme du bienheureux Pascal, votre Confesseur, d’un admirable et tendre amour pour les mystères sacrés de votre corps et de votre sang, accordez-nous, dans votre bonté, que nous méritions de retirer de ce banquet divin la même abondance de grâces qu’il y a trouvée.
$Qui vivis}
\festum{0518}{S. Venantii Martyris}{18 mai}
\saintindex{}
\saintindex{}
\oratio{Deus, qui hunc diem beáti Venántii Mártyris tui triúmpho consecrásti: exáudi preces pópuli tui et præsta: ut, qui ejus mérita venerámur, fídei constántiam imitémur.  \perdominumla}{Ô Dieu, qui avez consacré ce jour par le triomphe du bienheureux Venant, Votre Martyr, exaucez les prières de Votre peuple et faites qu’honorant ses mérites, nous imitions la constance de sa foi. \perdominumfr}
\festum{0518}{S. John I, pope and martyr}{18 mai}
\saintindex{}
\saintindex{}

\festum{0519}{S. Petri Celestini Papæ et Confessoris}{19 mai}
\saintindex{}
\saintindex{}
\oratio{Deus, qui beátum Petrum Cælestínum ad summi pontificátus ápicem sublimásti, quique illum humilitáti postpónere docuísti: concéde propítius; ut ejus exémplo cuncta mundi despícere, et ad promíssa humílibus prǽmia perveníre felíciter mereámur. \perdominumla}{Dieu, qui avez élevé le bienheureux Pierre Célestin à l’éminente dignité du souverain pontificat et qui lui avez appris à mettre l’humilité au-dessus de cette élévation, accordez-nous, dans Votre bonté, la grâce de mépriser, à son exemple, tous les biens de ce monde et de parvenir heureusement à la possession des récompenses promises à ceux qui sont humbles. \perdominumfr}
\oratio{@Commune/C4-1:Oratio:s/N./Petri Cœlestíni/}{}
\festum{0519}{S. Pudentianæ Virginis}{19 mai}
\saintindex{Pudentiánæ}
\saintindex{Pudentienne}
\festum{0520}{S. Bernardini Senensis Confessoris}{20 mai}
\saintindex{}
\saintindex{}
\oratio{Dómine Jesu, qui beáto Bernardíno Confessóri tuo exímium sancti nóminis tui amórem tribuísti: ejus, quǽsumus, méritis et intercessióne, spíritum nobis tuæ dilectiónis benígnus infúnde:
$Qui vivis}{Seigneur Jésus, qui avez accordé au bienheureux Bernardin, Votre Confesseur, un très ardent amour pour Votre saint Nom, nous Vous supplions, par ses mérites et son intercession, de daigner, dans Votre bonté, répandre en nous l’esprit de Votre charité.
$Qui vivis}
\festum{0521}{S. Christopher Magallanes and companions, martyrs}{21 mai}
\saintindex{}
\saintindex{}
\festum{0522}{S. Ritæ de Cascia}{22 mai}
\saintindex{}
\saintindex{}
\oratio{Deus, qui sanctæ Ritæ tantam grátiam conférre dignátus es, ut inimícos dilígeret et in corde ac fronte caritátis et passiónis tuæ signa portáret: da nobis, quǽsumus, ejus intercessióne et méritis; inimícis nostris sic párcere et passiónis tuæ dolóres contemplári, ut promíssa mítibus ac lugéntibus prǽmia consequámur:
$Qui vivis
}{}
\festum{0523}{S. Joannis Baptistæ de Rossi Confessoris}{23 mai}
\saintindex{}
\saintindex{}
\oratio{Deus, qui sanctum Joánnem Baptistam Confessorem tuum in evangelizandis pauperibus caritate et patientia decorasti: concede quǽsumus; ut cujus pia merita veneramur, virtutum quoque imitemur exempla. \perdominumla}{Dieu, qui avez orné saint Jean Baptiste, Votre Confesseur, de la charité et de la patience dans l’évangélisation des pauvres, accordez-nous, nous Vous le demandons, de vénérer ses pieux mérites et d’imiter les exemples de ses vertus. \perdominumfr}
\festum{0523}{S. John Baptist Rossi Confessor}{23 mai}
\saintindex{}
\saintindex{}
\oratio{Deus, qui sanctum Joánnem Baptistam Confessorem tuum in evangelizandis pauperibus caritate et patientia decorasti: concede qusesumus; ut cujus pia merita veneramur, virtutum quoque imitemur exempla. \perdominumla}{Dieu, qui avez orné saint Jean Baptiste, Votre Confesseur, de la charité et de la patience dans l’évangélisation des pauvres, accordez-nous, nous Vous le demandons, de vénérer ses pieux mérites et d’imiter les exemples de ses vertus. \perdominumfr}
\festum{0525}{S. Gregorii VII Papæ et Confessoris}{25 mai}
\saintindex{}
\saintindex{}
\oratio{Deus, in te sperántium fortitúdo, qui beátum Gregórium Confessórem tuum atque Pontíficem, pro tuénda Ecclésiæ libertáte, virtúte constántiæ roborásti: da nobis, ejus exémplo et intercessióne, ómnia adversántia fórtiter superáre. \perdominumla}{Ô Dieu, force de ceux qui espèrent en Vous, Vous qui avez fortifié le bienheureux Grégoire, Votre Confesseur et Pontife par la vertu de constance dans la défense de la liberté de l'Eglise, accordez-nous, par son exemple et son intercession, de triompher courageusement de toute adversité. \perdominumfr}
\festum{0525}{S. Urbanis Papæ et Martyris}{25 mai}
\saintindex{Urbáno}
\saintindex{Urbain}
\oratio{Da, quǽsumus, omnípotens Deus: ut qui beáti Urbáni Mártyris tui atque Pontíficis solémnia cólimus, ejus apud te intercessiónibus adjuvémur. \perdominumla}{Faites, nous Vous en prions, Dieu tout-puissant, que, nous qui célébrons la fête de Votre bienheureux Martyr et Pontife Urbain, nous recevions l’assistance de son intercession. \perdominumfr}
\oratio{@Commune/C2b:Oratio pro commemoratio
(sed rubrica 196)
@Commune/C2b}{Ô Dieu, sur le roc inébranlable des apôtres, vous avez fondé votre Eglise qui n’a plus a redouter la puissance des enfers ; faites qu’à la prière du bienheureux Urbain, votre Martyr et Souverain Pontife, elle reste ferme dans votre vérité et sous votre garde, elle soit en sûreté pour toujours. \perdominumfr
(sed rubrica 196 loco horum versum dicitur)
@Commune/C2b}
\festum{0526}{S. Philippi Neri Confessoris}{26 mai}
\saintindex{}
\saintindex{}
\oratio{Deus, qui beátum Philíppum Confessórem tuum Sanctórum tuórum glória sublimásti: concéde propítius; ut, cujus solemnitáte lætámur, ejus virtútum proficiámus exémplo. \perdominumla}{Dieu, qui avez élevé le bienheureux Philippe, Votre Confesseur, à la gloire de Vos Saints, accordez-nous, dans Votre miséricorde que, célébrant avec joie cette solennité, nous mettions à profit l’exemple de ses vertus. \perdominumfr}
\festum{0526}{S. Eleutherii Papæ et Martyris}{26 mai}
\saintindex{Eleuthérii}
\saintindex{Eleuthère}
\festum{0527}{S. Bedæ Venerabilis Confessoris et Ecclesiæ Doctoris}{27 mai}
\saintindex{Bedæ}
\saintindex{Bède}
\oratio{Deus, qui Ecclésiam tuam beáti Bedæ Confessóris tui atque Doctóris eruditióne claríficas: concéde propítius fámulis tuis; ejus semper illustrári sapiéntia et méritis adjuvári. \perdominumla}{Dieu, qui illustrez Votre Église par la science du bienheureux Bède, Votre Docteur et Confesseur, accordez dans Votre bonté à Vos serviteurs d’être toujours éclairés de sa sagesse et aidés de ses mérites. \perdominumfr}
\festum{0527}{S. Joannis Papæ et Martyris}{27 mai}
\saintindex{Joannis}
\saintindex{Jean}
\festum{0528}{S. Augustini Episcopi et Confessoris}{28 mai}
\saintindex{}
\saintindex{}
\oratio{Deus, qui Anglórum gentes, prædicatióne et miráculis beáti Augustíni Confessóris tui atque Pontíficis, veræ fídei luce illustráre dignátus es: concéde; ut, ipso interveniénte, errántium corda ad veritátis tuæ rédeant unitátem, et nos in tua simus voluntáte concórdes. \perdominumla}{Ô Dieu, qui, par la prédication et les miracles du bienheureux Augustin, Votre Confesseur et Pontife, avez daigné éclairer de la lumière de la vraie foi la nation anglaise, faites que, par son intercession, les cœurs égarés reviennent à l’unité de Votre vérité, et que nous soyons tous unis de cœur en Votre volonté. \perdominumfr}
\festum{0529}{S. Mariæ Magdalenæ de Pazzis Virginis}{29 mai}
\saintindex{}
\saintindex{}
\oratio{Deus, virginitátis amátor, qui beátam Maríam Magdalénam Vírginem, tuo amóre succénsam, cæléstibus donis decorásti: da; ut, quam festíva celebritáte venerámur, puritáte et caritáte imitémur. \perdominumla}{Dieu, qui aimez la virginité et qui avez orné des dons célestes Marie-Madeleine, cette vierge embrasée de Votre amour, donnez-nous d’imiter, dans sa pureté et sa charité, celle que nous vénérons en célébrant sa fête. \perdominumfr}
\festum{0530}{S. Felicis I Papæ et Martyris}{30 mai}
\saintindex{Felícis}
\saintindex{Félix}
\festum{0531}{}{31 mai}
\saintindex{}
\saintindex{}
\oratio{Concéde nobis, quǽsumus, Dómine: ut, qui solemnitátem beatæ Maríæ Vírginis Regínæ nostræ celebrámus; ejus muníti præsídio, pacem in præsénti et glóriam in futúro cónsequi mereámur. \perdominumla}{Nous célébrons la solennité de la Bienheureuse Vierge Marie, notre Reine : accordez-nous, Seigneur, de pouvoir, avec l'appui de son secours, obtenir la paix dans le présent et la gloire dans l'avenir. \perdominumfr}
\festum{0531}{S. Petronillæ Virginis}{31 mai}
\saintindex{Petroníllæ}
\saintindex{Petronille}
\festum{0601}{S. Angelæ Mericiæ Virginis}{1 juin}
\saintindex{}
\saintindex{}
\oratio{Deus, qui novum per beátam Angelam sacrárum Vírginum collégium in Ecclésia tua floréscere voluísti: da nobis, ejus intercessióne, angélicis móribus vívere; ut, terrénis ómnibus abdicátis, gáudiis pérfrui mereámur ætérnis. \perdominumla}{Ô Dieu, qui, par la bienheureuse Angèle, avez voulu qu’une nouvelle société de vierges saintes fleurît dans votre Eglise, faites-nous, par son intercession, la grâce de mener une vie angélique, afin que, renonçant à toutes les choses de la terre, nous méritions de jouir des joies éternelles. \perdominumfr}
\festum{0602}{Ss. Marcellini, Petri, atque Erasmi, Episcopi, Martyrum}{2 juin}
\saintindex{}
\saintindex{}
\oratio{Deus, qui nos ánnua beatórum Mártyrum tuórum Marcellíni, Petri atque Erásmi solemnitáte lætíficas: præsta, quǽsumus; ut, quorum gaudémus méritis, accendámur exémplis. \perdominumla}{Ô Dieu, qui nous faites trouver un sujet de joie dans la solennité de Vos bienheureux Martyrs Marcellin, Pierre et Erasme, accordez-nous, s’il Vous plaît, la grâce d’être enflammés d’ardeur par les exemples de ceux dont les mérites nous réjouissent.  \perdominumfr}
\festum{0603}{Ss. Caroli Lwanga, Matthiæ Murumba et Sociorum Martyrum}{3 juin}
\saintindex{}
\saintindex{}
\oratio{Omnípotens et miséricors Deus, qui univérsum mundum lúmine veritátis illústras, et per gloriósos Mártyres Matthíam, Cárolum et Sócios fídei prodígia in Africæ regiónibus renováre dignátus es: largíre supplícibus tuis; ut, ipsis intercedéntibus, in nobis grátiæ tuæ dona custódias, et inter gentes omnes mérito et número pópulus tibi sérviens augeátur. \perdominumla}{}
\festum{0604}{S. Francisci Caracciolo Confessoris}{4 juin}
\saintindex{}
\saintindex{}
\oratio{Deus, qui beátum Francíscum, novi órdinis institutórem, orándi stúdio et pœniténtiæ amóre decorásti: da fámulis tuis in ejus imitatióne ita profícere; ut, semper orántes et corpus in servitútem redigéntes, ad cæléstem glóriam perveníre mereántur. \perdominumla}{Ô Dieu, qui avez suscité le bienheureux François pour être le fondateur d’un nouvel Ordre et l’avez admirablement doué de zèle pour la prière et d’amour de la pénitence, accordez à Vos serviteurs de profiter si bien de ses exemples que, s’appliquant toujours à prier et à réduire leur corps en servitude, ils méritent de parvenir à la gloire céleste. \perdominumfr}
\festum{0605}{S. Bonifatii Episcopi et Martyris}{5 juin}
\saintindex{}
\saintindex{}
\oratio{Deus, qui multitúdinem populórum, beáti Bonifátii Mártyris tui atque Pontíficis zelo, ad agnitiónem tui nóminis vocáre dignátus es: concéde propítius; ut, cujus solémnia cólimus, étiam patrocínia sentiámus. \perdominumla}{Ô Dieu, qui avez daigné appeler une multitude de peuples à la connaissance de Votre nom par le zèle du bienheureux Boniface, Votre Martyr et Pontife, accordez-nous, dans Votre bonté que, célébrant sa fête, nous ressentions les effets de sa protection. \perdominumfr}
\festum{0606}{S. Norberti Episcopi et Confessoris}{6 juin}
\saintindex{}
\saintindex{}
\oratio{Deus, qui beátum Norbértum Confessórem tuum atque Pontíficem verbi tui præcónem exímium effecísti, et per eum Ecclésiam tuam nova prole fecundásti: præsta, quǽsumus; ut, ejúsdem suffragántibus méritis, quod ore simul et ópere dócuit, te adjuvánte, exercére valeámus. \perdominumla.}{Dieu, Vous avez fait du bienheureux Norbert, Votre Confesseur et Pontife, un excellent prédicateur de Votre parole et Vous avez donné par lui à Votre Église une nouvelle famille : faites, nous Vous en supplions, qu’aidés de ses mérites, nous puissions, grâce à Votre secours, mettre en pratique ce qu’il a enseigné par ses paroles et par ses œuvres. \perdominumfr.}
\festum{0609}{Ss. Primi et Feliciani Martyrum}{9 juin}
\saintindex{}
\saintindex{}
\oratio{Fac nos, quǽsumus, Dómine, sanctórum Mártyrum tuórum Primi et Feliciáni semper festa sectári: quorum suffrágiis protectiónis tuæ dona sentiámus. \perdominumla}{Nous Vous en supplions, Seigneur, faites que nous célébrions toujours fidèlement la fête de Vos saints Martyrs Prime et Félicien, afin que, grâce à leur intercession, nous éprouvions les bienfaits de Votre protection. \perdominumfr}
\festum{0610}{S. Margaritæ Reginæ Viduæ}{10 juin}
\saintindex{}
\saintindex{}
\oratio{Deus, qui beátam Margarítam regínam exímia in páuperes caritáte mirábilem effecísti: da; ut, ejus intercessióne et exémplo, tua in córdibus nostris cáritas júgiter augeátur. \perdominumla}{Dieu, Vous avez rendu admirable la bienheureuse reine Marguerite en lui inspirant une extrême charité pour les pauvres : faites que, par son intercession et à son exemple, Votre charité croisse continuellement dans nos cœurs. \perdominumfr}
\festum{0611}{S. Barnabæ Apostoli}{11 juin}
\saintindex{}
\saintindex{}
\oratio{Deus, qui nos beáti Bárnabæ Apóstoli tui méritis et intercessióne lætíficas: concéde propítius; ut, qui tua per eum benefícia póscimus, dono tuæ grátiæ consequámur. \perdominumla}{Dieu, Vous nous donnez un motif de joie dans les mérites et l’intercession du bienheureux Barnabé, Votre Apôtre : accordez-nous, avec bonté, qu’en recourant à cette intercession pour solliciter Vos bienfaits, nous les obtenions au moyen de Votre grâce. \perdominumfr}
\festum{0612}{S. Joannis a S. Facundo Confessoris}{12 juin}
\saintindex{}
\saintindex{}
\oratio{Deus, auctor pacis et amátor caritátis, qui beátum Joánnem Confessórem tuum mirífica dissidéntes componéndi grátia decorásti: ejus méritis et intercessióne concéde; ut, in tua caritáte firmáti, nullis a te tentatiónibus separémur. \perdominumla}{Dieu, auteur de la paix et amant de la charité, Vous avez orné le bienheureux Jean, Votre Confesseur, d’un merveilleux don du ciel pour apaiser les différends : accordez-nous, par ses mérites et son intercession, d’être tellement affermis dans Votre amour, que nous ne soyons plus séparés de Vous par aucune tentation. \perdominumfr}
\festum{0612}{Ss. Basilidis, Cyrini, Naboris and Nazarii Martyrum}{12 juin}
\saintindex{}
\saintindex{}
\oratio{Sanctórum Mártyrum tuórum Basílidis, Cyríni, Náboris atque Nazárii, quǽsumus, Dómine, natalítia nobis votíva respléndeant: et, quod illis cóntulit excelléntia sempitérna, frúctibus nostræ devotiónis accréscat. \perdominumla}{Faites, s’il Vous plaît, Seigneur, que l’anniversaire de la naissance au ciel de Vos saints Martyrs Basilide, Cyrin, Nabor et Nazaire répande la lumière dans nos âmes et que les biens spirituels dont le bonheur éternel leur a procuré la plénitude, s’accroissent en nous grâce aux fruits de notre dévotion. \perdominumfr}